\documentclass[11pt,a4paper]{article}
\usepackage[T1]{fontenc}\usepackage{lmodern,amsmath,amssymb,amsthm,mathtools,booktabs,longtable,microtype}
\usepackage[margin=25mm]{geometry}\usepackage[colorlinks=true,linkcolor=blue,urlcolor=blue,citecolor=blue]{hyperref}
\setlength{\emergencystretch}{3em}
\newtheorem{theorem}{Theorem}[section]\newtheorem{lemma}[theorem]{Lemma}\newtheorem{proposition}[theorem]{Proposition}\newtheorem{corollary}[theorem]{Corollary}
\newcommand{\Q}{\mathbb Q}\newcommand{\Z}{\mathbb Z}\newcommand{\Nm}{\operatorname{N}}
\title{Quartics in 196 and 258 variables by elimination after monomial rescaling}
\author{}\date{15 November 2026}
\begin{document}\maketitle
\begin{abstract}
Two exact changes of coordinates remove stored fourth powers from existing
quadratic systems without producing cubic equations. The 198-variable family
supplied in the Conant--ChatGPT research archive becomes a 196-variable
family with 182 quadratic equations. The fully specified 260-variable family
becomes a 258-variable family with 230 quadratic equations. These changes
preserve the entire rational input projection for every fixed coefficient
parameter, rather than just a chosen completeness set. The latter family
therefore retains the certified defining exponent $N=9682$. Its largest
coefficient has 139819 binary digits. The former defining threshold remains
ineffective. We prove the nonzero conditions before division and specify all
maps, equations, and coefficients by exact finite data with separate checking
algorithms. No new arithmetic hypothesis is introduced by these transformations;
the parent projection theorems retain their stated arithmetic premises.
\end{abstract}

\section{Statements and scope}
A total variable count includes the input $X$. A coefficient parameter is
fixed and is not a quantified variable. Write $\pi_X(V(\Q))$ for the set of
rational inputs of a rational algebraic zero set $V$. All statements about
integer definitions below are conditional on the precise arithmetic inputs
of the indicated parent constructions. The algebraic equivalences themselves
hold without those arithmetic inputs.

The first parent is the unbalanced table supplied in
\texttt{H10-quartic198.zip}, namely
\begin{center}\small\path{circuits/inverse198/quadratic.json}.\end{center}
Its SHA-256 is
\begin{center}\small\texttt{c55d7eefe41f62fee88b467e3ea8f6be8ba1e66f6a94b30eb74b39f6f6e7efad}.\end{center}
The archive attributes its new analysis and code to ChatGPT in collaboration
with Gabriel Conant's research request. Its 198-variable theorem builds on
the project's 204-variable construction, not on the unpublished erroneous
scalar substitution previously described as v32. The original editable sources and data are
preserved without changes under \path{sources/h10_quartic198}; its own
attribution, theorem, qualifications, and source bindings are retained.
The original archived PDF is identified by its byte length and hash in
\path{PDF_PROVENANCE.md}, rather than redistributed in this TeX-first release.

The second parent is the complete 260-variable table of \cite{Explicit}, at
repository commit \texttt{e950bd2eb46ffb6aaf77624cc9b3d0b880fd921f}.
Both parents use $A=149^{2N+1}$, with integer $N\ge1$.
Set $C=149^N$, so $A=149C^2$, and put
\begin{align*}
 K_{196}&=11903780636513165769654680680436132271178681610029854792,\\
 K_{\rm fixed}&=484991541504995159662206970923766531374664106904986910720.
\end{align*}

\begin{theorem}\label{thm:results}
There are explicit families of integer polynomials $P_N$ and $Q_N$ with
\[
\begin{array}{c|r|r|r|r}
 &\text{total variables}&\text{quadratic summands}&\text{degree}&\text{monomials}\\\hline
P_N&196&182&4&2943\\
Q_N&258&230&4&3705
\end{array}
\]
for every $N\ge1$. Both have integer content one, and
\[
 H(P_N)=K_{196}149^{2N},\qquad
 H(Q_N)=\max\{K_{\rm fixed},9792^2 149^{2N}\}.
\]
For every fixed $N\ge1$, $P_N$ has the same rational input projection as
the supplied 198-variable parent, and $Q_N$ has the same rational input
projection as the fully specified 260-variable parent. Consequently:
\begin{enumerate}
\item Under the inherited arithmetic premises of \cite{Archive198}, every
sufficiently large fixed $N$ makes $P_N$ define $\Z$ over $\Q$. The sufficient
threshold is the same as the parent's and remains ineffective.
\item Under the inherited arithmetic premises of \cite{Explicit}, every
$N\ge9682$ makes $Q_N$ define $\Z$. At $N=9682$ its exact height is
$9792^2 149^{19364}$, whose binary length is 139819.
\end{enumerate}
The coefficient crossover in the formula for $H(Q_N)$ is $N=12$; it is not a
sufficient defining threshold. No 233-variable theorem is asserted here.
\end{theorem}

The inherited premises include the specified elliptic group and index
statements over $\Q(\sqrt2)$, the appropriate contact-height/modular inputs,
the pole obstruction and prepared pole realization, and fixed-data
simultaneous prime selection. The 196-variable parent permits half-integral
elliptic indices and finite torsion differences as stated in its proof.
We do not narrow that domain. The 258-variable theorem uses exactly the
effective threshold of the 260-variable parent. Neither statement adds $abc$.
The separate 78-variable family \cite{ABC78} still uses ordinary rational
$abc$ and its own stated group, prime, and simultaneous-twist inputs.

\section{Elimination on a locus where specified coordinates are nonzero}
\begin{theorem}\label{thm:general}
Let a rational system $f_r=0$ have distinguished input $X$, retained
coordinates $u_1,\ldots,u_k$, other coordinates $x_i$, and some coordinates
$t_j$ to be deleted. Suppose specified retained equations imply
$u_1\cdots u_k\ne0$. Suppose that, on this locus, the defining equations
for the $t_j$ uniquely give monomials in the $u_\ell$ with nonzero rational
coefficients. Make a diagonal substitution
\[
 x_i=c_i u_1^{w_{i1}}\cdots u_k^{w_{ik}}y_i,
 \qquad c_i\in\Q^\times,\quad w_{i\ell}\in\Z,
\]
fixing $X$ and the $u_\ell$, and substitute the formulas for the $t_j$.
For every remaining row suppose its substituted Laurent polynomial is
\[
 f_r(\text{substitution})=d_r\boldsymbol u^{\boldsymbol a_r}g_r,
 \qquad d_r\in\Q^\times,
\]
where $g_r$ is an ordinary polynomial of degree at most two. Suppose also
that the new equations imply the same required nonvanishing and retained
power identities. Then the two rational zero sets are in bijection by
these substitutions, with inverse given by inverse diagonal scaling and
the unique formulas for the deleted coordinates. In particular their
rational input projections agree. Clearing rational coefficients of each
$g_r$ to primitive integers and taking their sum of squares gives a
polynomial of degree at most four with that same projection.
\end{theorem}
\begin{proof}
Every old solution has nonzero $u_\ell$, so inverse scaling defines a
new tuple. The nonzero multipliers in the displayed identities show that
all new equations vanish. Conversely, on a new solution all $u_\ell$ are
nonzero by hypothesis. Reconstruct the $x_i$ and $t_j$ by the formulas.
Deleted defining equations hold identically; every other old row vanishes
by its identity with $g_r$. The two maps undo one another. Since the input
is fixed, their projections agree. Over $\Q$ a sum of rational squares
vanishes exactly when every summand vanishes, proving the last assertion.
\end{proof}

No density argument is involved: this is an all-solution statement.
Nonvanishing must follow from the new equations, not just from the old
system. This is why the retained inverse equations are checked separately.

\subsection{A finite criterion for keeping the rows quadratic}
When the retained power coordinates have exponent vectors in a finite set
$\mathcal P\subset\Z^k$, let $\mathcal R_c$ be the vectors obtainable as
sums of at most $c$ elements of $\mathcal P$, including the empty sum.
Collect all cancellations in a substituted row and write its terms as
$a_{m,e}y^m\boldsymbol u^e$. For a fixed row multiplier
$\boldsymbol u^{-s}$, each such term has a quadratic representation by
retained power coordinates whenever
\[
 |m|\le2,\qquad e-s\in\mathcal R_{2-|m|}.
\]
Thus the finite intersection of the sets $e-\mathcal R_{2-|m|}$, over its
nonzero terms, supplies all possible row shifts for this particular
monomial representation. Multiplication by each selected power monomial
is verified using its defining relation. This criterion is not claimed
necessary among arbitrary consequences of the whole ideal, and an
unsuccessful search gives no lower bound on variable counts.

\section{The 196-variable construction}
\subsection{The nonzero scalar normalizer}
In the 198-variable parent let $R$ be the scalar normalizer, $J$ the
rescaled contact product, $u$ the input-square coordinate, and $v$ the
first scalar-product coordinate. The following four literal rows are
retained without scaling:
\begin{align*}
 u&=X^2,\\
 4v&=6255001+255102u+2601u^2,\\
 0&=14315517724+291790992u+15531948uv+767410608v-512v^2+RJ,\\
 0&=1+22201RJ-AJ.
\end{align*}
Substituting the first two into the third gives exactly
\[
 RJ=(51X^2+103)(51X^2+2503)(51X^2+2504)(1632X^2+80471)=r(X)>0.
\]
Hence $R,J\ne0$, and the fourth row gives
\[
 R=\frac{Ar(X)}{1+22201r(X)}>0.
\]
The retained square row gives $R_2=R^2$. This argument is an elementary
identity and positivity proof valid for every rational input, before
any source arithmetic theorem is invoked.

\subsection{The detector coordinates and substitutions}
For detector $j=0,1$, write its retained inverse weight as $w_j$ and its
contact weight as $q_j$. Retain literally
\[
 q_jw_j=1,\qquad w_{j,2}=w_j^2.
\]
The two deleted fourth-power coordinates have physical values
\[
 w_{0,4}=w_0^4,\qquad w_{1,4}=w_1^4/R.
\]
The factor $R^{-1}$ in the second formula is part of the parent storage
convention, not an additional assumption. The corresponding two defining
rows become identities and are removed.

Within detector $j$, old values of the following named coordinates equal
the indicated factor times their new values:
\[
\begin{array}{c|c}
\text{coordinate in the source table}&\text{factor}\\\hline
\text{normalized.s},\ \text{factored.v}&w_j^{-1}\\
\text{normalized.p}&w_j\\
\text{normalized.p2.quadratic}&w_j^2\\
\text{actualW.0},\ \text{actualW.1},\ \text{normOver2w0},\ \text{B}&w_j^{-2}\\
\text{amplified.det}j\text{.factor}k\text{.}c\ (0\le k\le4)&R^j w_j^{-2}\\
\text{amplified.det}j\text{.product1.}c&w_j^{-2}\\
\text{amplified.det}j\text{.product}k\text{.}c\ (k=2,3)&R^{-j}w_j^{-2}
\end{array}
\]
Here $c=0,1$ denotes the rational component. All unspecified coordinates
are fixed. Put $H=31581660^{14}$ and $K=82479248H$. A simultaneous
constant scaling multiplies old scalar coordinates, in terms of their
new values, by
\[
\begin{array}{c|c}
\text{scalar.W2.quadratic.rescaled}&H^{-1}\\
\text{scalar.b.quadratic}&-K^{-1}\\
\text{normalized\_scalar.scaled\_e}&-K\\
\text{normalized\_scalar.s},\ \text{normalized\_scalar.p},\
 \text{normalized\_scalar.factored.q}&K\\
\text{normalized\_scalar.p2.quadratic},\ \text{normalized\_scalar.factored.q2}&K^2\\
\text{normalized\_scalar.factored.v}&-1
\end{array}
\]
These constants are nonzero and none of the nine nonzero-guard and
power rows is changed. There are $198-2=196$ coordinates and
$184-2=182$ retained rows.

The exact coordinate indices, row multipliers, primitive clearing factors,
and final rows are supplied in
\path{computation/quartic/laurent_recovery_v38/config196.json} and
\path{family196/certificate.json} relative to that directory. For each
row its Laurent expansion and target expansion coincide coefficient by
coefficient. The certificate lists exactly one identity for each retained
row; every source row is either listed or one of the two deleted
identities. This is a finite proof of the required row identities,
reproducible by the independent checker in Section~\ref{sec:verification}.
The coordinate table, not a generic existence assertion, specifies every
coefficient of the defining polynomial.

All three base units $R,w_0,w_1$ are nonzero on both systems by the
literal equations proved above. Theorem~\ref{thm:general} therefore
applies on the entire rational zero set, including every original
identity, tangency, or torsion stratum. It gives an input-preserving
bijection for each $N\ge1$.

\section{The fully specified 258-variable construction}
The first step starts with \path{compressed_reference_v27/family260/quadratic.json}.
Within its detector zero, $q_0w_0=1$ and $w_{0,2}=w_0^2$ are literal
retained rows. Eliminate $w_{0,4}=w_0^4$. Rescale the remaining coordinates
by the integer powers of $w_0$ in \path{config259.json}. There are 259
coordinates and 231 quadrics. In the second step, the detector-one
relations are still literally $q_1w_1=1$, $w_{1,2}=w_1^2$.
Eliminate $w_{1,4}=w_1^4$ and use the powers in \path{config258.json}.
This gives 258 coordinates and 230 quadrics.

For completeness, Appendix~\ref{app:weights} prints every nonzero
coordinate exponent in these two steps. It is important to perform the
steps in order: the second table uses the coordinate names after the
first step. Their finite row certificates, generated and checked
separately, give all common row powers and primitive integer factors.

At the second step $q_1w_1=1$ remains unchanged. The other two inverse
relations in the final table are
\[
 q_0w_0=w_1,\qquad q_2w_2=w_1^2.
\]
Thus all inverse weights are still nonzero, without a new coordinate.
The two square relations used to reconstruct the deleted powers also
remain valid. Apply Theorem~\ref{thm:general} first to the 260-to-259 map,
then to the 259-to-258 map. At each step the specific unit used for the
change of coordinates is excluded from zero by a retained literal inverse
row. The composition is consequently a bijection on rational solutions.
No elliptic index, point chart, local order, or comparison argument is
changed as a physical quantity. In particular the effective parent
threshold $9682$ is inherited unchanged; it is not recomputed from a
new height inequality.

\section{Coefficient balancing and exact heights}
For a fixed nonzero $C$, a further diagonal substitution
$x_i=c_i C^{v_i}y_i$ is invertible. Substitute $A=149C^2$, divide each
row by its least power of $C$, and clear rational coefficients to
primitive integers. Neither operation changes its rational zero set.
The complete exponent vectors and row multipliers are in
\path{family196/balancing.json} and \path{family258/balancing.json}.
The 258-variable balancing uses only three nonzero $C$ weights: the
coordinates named \texttt{scalar.r\_power\_2},
\texttt{scalar.r\_power\_4}, and
\texttt{scalar.forced\_positive\_parameter} each have weight one.
Their old semantic names do not assert that their new stored values
are literal powers before reconstruction. The scalar constant scalings
are as above, with \texttt{scalar.W2.quadratic} in place of its
198-parent name.

For each coordinate monomial $m$, let $p_m(C)$ be its integer coefficient
after expanding the sum of squares and collecting equal monomials.
Every $p_m$ has degree at most two. The coefficient of $X^4$ is exactly
one, and all coordinates occur. Thus the degree is exactly four and
the integer content is one for every fixed integer $C=149^N$.

For the 196-variable system the independent coefficient reconstruction
verifies, for every $m$,
\[
 \sum_a |[C^a]p_m|149^a\le K_{196}149^2.
\]
Since $a\le2$, division by $C^2$ shows the same bound at all $C\ge149$.
A pure coefficient $K_{196}C^2$ occurs, proving equality for the height.
For each profile the highest-power term at $C=149$ strictly dominates
the sum of the absolute values of the lower terms; it continues to do
so for every larger $C$. Hence no coefficient cancels and the 2943
monomial count holds for every $N\ge1$.

For the 258-variable system the largest fixed coefficient is
$K_{\rm fixed}$ and is attained by a constant profile. A pure coefficient
$9792^2 C^2$ is also attained. Exact integer comparisons give
\[
 9792^2 149^{22}<K_{\rm fixed}<9792^2 149^{24}.
\]
Every profile satisfies the absolute-coefficient bound by $K_{\rm fixed}$
at $C=149^{11}$, and the bound by $9792^2C^2$ at $C=149^{12}$.
The former bound extends downwards to $C=149$, the latter upwards to
all larger $C$, proving the maximum formula of Theorem~\ref{thm:results}.
The same leading-term domination test excludes zeros for all profiles
at $C\ge149$. There are 3705 monomials. Finally the integer comparison
\[
 2^{139818}\le9792^2 149^{19364}<2^{139819}
\]
proves the stated binary length exactly; no floating-point logarithm
or decimal rounding is used.

\section{Independent reconstruction and limitations}\label{sec:verification}
The producer uses sparse Laurent exponent arithmetic. The checker does not
import it: it separately expands the two sides of each row identity,
verifies all necessary unit and power relations, checks row and coordinate
coverage, and expands each sum of squares using unordered term pairs.
It checks every coefficient, not selected evaluations. Each complete
producer is rerun in a temporary directory and the regenerated files are
compared byte for byte. There are five deliberate corruption tests for
the 196-variable construction and eight for the two explicit reductions.
Altered coefficients, shifts, unit guards, or powers are rejected.

The supplied 198-variable archive has its own thirteen-command verifier;
that verifier is run separately, in a temporary copy. The source archive's
infinite arithmetic inputs are not proved by any of these finite checks.
The 78-variable mixed-height construction is reconstructed by its separate
checker and retains ordinary rational $abc$. The general counts 25 and
13 are unchanged. Their finite-algebra presentations already eliminate
stored products as finite choices, so the present deletion of auxiliary
power coordinates alone does not remove one of their retained point or
pole parameters.

The two reattached v36 recovery archives contain neither a 233-variable
manuscript nor its coefficient tables. Their own receipt states that
recovery was not closed. That earlier conversational number, its asserted
exponent 10220, and its asserted torsion-height estimate are not used as
inputs to this work and are not certified here. The source inventory and
hashes are recorded separately in the recovery audit.

\appendix
\section{All nonzero exponents for the two explicit deletions}\label{app:weights}
In each table an exponent $v$ means $x_{\rm old}=w_j^v x_{\rm new}$;
all omitted coordinates are fixed. Stored powers are interpreted by the
separate retained power equations, as specified in the proof.
\subsection*{260 to 259, unit $w_0$}
\small
\begin{longtable}{p{0.80\linewidth}r}
\toprule Coordinate name&Exponent\\\midrule\endhead
\texttt{detector0.normalized.s} & -1\\
\texttt{detector0.actualW.0} & -2\\
\texttt{detector0.actualW.1} & -2\\
\texttt{detector0.normOver2w0} & -2\\
\texttt{detector0.B} & -2\\
\texttt{detector0.normalized.p2.quadratic} & 2\\
\texttt{detector0.factored.v} & -1\\
\texttt{amplified.det0.factor0.0} & -2\\
\texttt{amplified.det0.factor0.1} & -2\\
\texttt{amplified.det0.factor1.0} & -2\\
\texttt{amplified.det0.factor1.1} & -2\\
\texttt{amplified.det0.factor2.0} & -2\\
\texttt{amplified.det0.factor2.1} & -2\\
\texttt{amplified.det0.factor3.0} & -2\\
\texttt{amplified.det0.factor3.1} & -2\\
\texttt{amplified.det0.product1.0} & -2\\
\texttt{amplified.det0.product1.1} & -2\\
\texttt{amplified.det0.product2.0} & -2\\
\texttt{amplified.det0.product2.1} & -2\\
\texttt{comparison.C2.fourthPower.0} & 1\\
\texttt{comparison.C2.fourthPower.1} & 1\\
\texttt{comparison.G.fourthPower.0} & 1\\
\texttt{comparison.G.fourthPower.1} & 1\\
\texttt{comparison.C3.fourthPower.0} & 1\\
\texttt{comparison.C3.fourthPower.1} & 1\\
\texttt{comparison.C4.fourthPower.0} & 1\\
\texttt{comparison.C4.fourthPower.1} & 1\\
\texttt{detector0.weightFifthPower} & -1\\
\texttt{detector1.weightFifthPower} & -1\\
\texttt{detector2.weightFifthPower} & -1\\
\bottomrule\end{longtable}
\normalsize
\subsection*{259 to 258, unit $w_1$}
\small
\begin{longtable}{p{0.80\linewidth}r}
\toprule Coordinate name&Exponent\\\midrule\endhead
\texttt{detector1.normalized.s} & -1\\
\texttt{Lucas.F0.quadratic} & 1\\
\texttt{Lucas.F1.quadratic} & 1\\
\texttt{testSquare.40.1.0.quadratic} & 1\\
\texttt{testSquare.40.1.1.quadratic} & 1\\
\texttt{testSquare.40.2.0.quadratic} & 1\\
\texttt{testSquare.40.2.1.quadratic} & 1\\
\texttt{testSquare.40.3.0.quadratic} & 1\\
\texttt{testSquare.40.3.1.quadratic} & 1\\
\texttt{testSquare.40.4.0.quadratic} & 1\\
\texttt{testSquare.40.4.1.quadratic} & 1\\
\texttt{detector0.q1} & -1\\
\texttt{detector1.actualW.0} & -2\\
\texttt{detector1.actualW.1} & -2\\
\texttt{detector1.normOver2w0} & -2\\
\texttt{detector1.B} & -2\\
\texttt{detector1.normalized.p2.quadratic} & 2\\
\texttt{detector1.factored.v} & -1\\
\texttt{amplified.det1.factor0.0} & -2\\
\texttt{amplified.det1.factor0.1} & -2\\
\texttt{amplified.det1.factor1.0} & -2\\
\texttt{amplified.det1.factor1.1} & -2\\
\texttt{amplified.det1.factor2.0} & -2\\
\texttt{amplified.det1.factor2.1} & -2\\
\texttt{amplified.det1.factor3.0} & -2\\
\texttt{amplified.det1.factor3.1} & -2\\
\texttt{amplified.det1.product1.0} & -2\\
\texttt{amplified.det1.product1.1} & -2\\
\texttt{amplified.det1.product2.0} & -2\\
\texttt{amplified.det1.product2.1} & -2\\
\texttt{detector2.q1} & -2\\
\texttt{detector2.actualW.0} & -1\\
\texttt{detector2.actualW.1} & -1\\
\texttt{detector2.normOver2w0} & -1\\
\texttt{comparison.C2.fourthPower.0} & 2\\
\texttt{comparison.C2.fourthPower.1} & 2\\
\texttt{comparison.G.fourthPower.0} & 2\\
\texttt{comparison.G.fourthPower.1} & 2\\
\texttt{comparison.C3.fourthPower.0} & 2\\
\texttt{comparison.C3.fourthPower.1} & 2\\
\texttt{comparison.C4.fourthPower.0} & 2\\
\texttt{comparison.C4.fourthPower.1} & 2\\
\texttt{detector0.weightFifthPower} & -2\\
\texttt{detector1.weightFifthPower} & -2\\
\texttt{detector2.weightFifthPower} & -2\\
\bottomrule\end{longtable}
\normalsize


\begin{thebibliography}{9}
\bibitem{Archive198} Gabriel Conant research request with ChatGPT,
\emph{A quartic in 198 variables}, 14 November 2026, supplied archive
\texttt{H10-quartic198.zip}; original proof and attribution retained in
\path{sources/h10_quartic198/arithmetic/inverse_vieta_198.tex} and its README.
\bibitem{Explicit} \emph{A fully specified defining quartic in 260 variables},
project source, 14 November 2026,
\path{papers/quartic/explicit260_20261114_v27.tex}.
\bibitem{ABC78} \emph{A quartic in 78 variables from a mixed-power height estimate},
project source, 15 November 2026,
\path{papers/quartic_existence/quartic_abc78_20261115_v34.tex}.
\end{thebibliography}
\end{document}
